\message{ !name(math.tex)}\documentclass[12pt]{article}
\usepackage{xeCJK}
\usepackage[utf8]{inputenc}
\usepackage{amsmath, amssymb, amsthm, mathtools, subfiles}

% Tree diagram
\usepackage{tikz}
\usetikzlibrary{trees}
\usepackage{verbatim}

% Drawing
\usepackage{pgf,pgfplots}
\pgfplotsset{compat=1.15}
\usepackage{mathrsfs}
\usetikzlibrary{arrows}
\pagestyle{empty}
\definecolor{rvwvcq}{rgb}{0.08235294117647059,0.396078431372549,0.7529411764705882}

% Hyperlink ref
% Be careful, it has to be the last package to be imported
\usepackage{hyperref}
% PDF-specific options
\hypersetup{
    colorlinks=true,
    linkcolor=blue,
    filecolor=magenta,
    urlcolor=cyan,
    pdftitle={Math},
    bookmarks=true,
}

% simple way to define
\usepackage{geometry}
\geometry{
  a4paper,
  total={170mm,257mm},
  left=20mm,
  top=20mm,
}


\setlength{\parindent}{0em} % paragragh indent
\setlength{\parskip}{0.6em}  % or use \bigskip \medskip \smallskip in content
\linespread{1.25} % line spacing

\title{MATH}
% \date{2019}
% \author{zhb}

\begin{document}

\message{ !name(math.tex) !offset(155) }
$
\left \{
    \begin{array}{l}
      A_1x+B_1y+C_1z=0\\
      A_2x+B_2y+C_2z=0
    \end{array}
  \right.
$\\
$\Rightarrow
\lambda(A_1x+B_1y+C_1z) + \mu(A_2x+B_2y+C_2z)=0$, $\lambda$ and $\mu$ 不同时为0\\
$\Rightarrow A_1x+B_1y+C_1z+\lambda(A_2x+B_2y+C_2z)=0$, same as above? sure...

\message{ !name(math.tex) !offset(257) }

\end{document}
%%% Local Variables:
%%% mode: latex
%%% TeX-PDF-mode: t
%%% TeX-engine: xetex
%%% TeX-master: t
%%% End:
